\uuid{bf1i}
\exo7id{1633}
\titre{exo7 1633}
\auteur{barraud}
\organisation{exo7}
\datecreate{2003-09-01}
\isIndication{false}
\isCorrection{false}
\chapitre{Réduction d'endomorphisme, polynôme annulateur}
\sousChapitre{Diagonalisation}
\cpge{mp-mpi.reduction-endomorphismes.b, pc.reduction.a, psi.reduction-endomorphismes.a, pt.reduction-endomorphismes.a, mp-mpi.reduction-endomorphismes.c, pc.reduction.b, psi.reduction-endomorphismes.b, pt.reduction-endomorphismes.b, mp-mpi.reduction-endomorphismes.d, pc.reduction.c, psi.reduction-endomorphismes.c, pt.reduction-endomorphismes.c, mp-mpi.reduction-endomorphismes.e, pc.reduction.e, psi.reduction-endomorphismes.e, pt.reduction-endomorphismes.d, bcpst2.valeurs-propres-vecteurs-propres.a, bcpst2.valeurs-propres-vecteurs-propres.b}
\module{Algèbre}
\niveau{L2}
\difficulte{}

\contenu{
\question{
Les matrices suivantes sont-elles diagonalisables, triangularisables,
sur $\R$ et sur $\C$ ?

Lorsqu'elles sont diagonalisables, donner une matrice diagonale semblable.

$$
 A=
 \begin{pmatrix}
   3 & 2 &-1 &-2 \\
   1 & 3 &-1 &-1 \\
   2 & 2 & 0 &-2 \\
   1 & 2 &-1 & 0
 \end{pmatrix}
\qquad
 B=  
 \begin{pmatrix}
   1 & 2 &-1 & 1 \\
   1 & 1 & 1 &-1 \\
   3 &-4 & 5 &-3 \\
   0 & 0 & 0 & 2
 \end{pmatrix}
\qquad
 C=  
 \begin{pmatrix}
  -1 & 0 & 1 \\
   1 &-1 & 0 \\
  -4 & 2 & 2 
 \end{pmatrix}
$$

Réduire explicitement $A$ et $C$.
}
}
